\documentclass[aspectratio=169]{beamer}
\usetheme{Madrid}
\usecolortheme{default}
\usepackage{booktabs}
\usepackage{graphicx}
\usepackage{siunitx}
\sisetup{round-mode=places,round-precision=3}

\title{ALEPH Thrust Unfolding Study\\Iteration Scan and $\chi^2/\mathrm{ndf}$ Checks}
\author{Graduate Student Study Note}
\date{\today}

\begin{document}

\begin{frame}
  \titlepage
\end{frame}

\begin{frame}{Recent Comments Summary}
\begin{itemize}
  \item We agreed the convergence decision should include iteration-to-iteration stability, not only closure/refold metrics.
  \item Added metric: $\chi^2/\mathrm{ndf}(\mathrm{unfolded}_{N},\mathrm{unfolded}_{N-1})$ to quantify incremental change.
  \item Observed smooth monotonic decrease of incremental change with iteration.
  \item Practical convergence threshold crossing occurs at $N_{\mathrm{iter}}=7$.
  \item This is consistent with the previous compromise point from refold + toy closure diagnostics.
\end{itemize}
\end{frame}

\begin{frame}{Study Setup}
\begin{itemize}
  \item Objective: choose Bayesian unfolding iterations using $\chi^2/\mathrm{ndf}$ diagnostics.
  \item Config: \texttt{Unfolding/unfold\_config.json}, scan iterations $N_{\mathrm{iter}}=1\dots 20$.
  \item Inputs: 1994 data reco, MC reco, MC gen (same files as nominal unfolding).
  \item Diagnostics used:
  \begin{itemize}
    \item Refold test: $\chi^2/\mathrm{ndf}(\text{data reco},\,\text{refolded unfolded})$.
    \item Toy closure: Poisson fluctuations of reco expectation from response$\times$truth, then unfold each toy and compare to MC truth.
  \end{itemize}
  \item Toy statistics: 60 toys per iteration, seed 12345.
\end{itemize}
\end{frame}

\begin{frame}{Verification of 4--5 Iteration Claim}
\begin{itemize}
  \item Cross-check performed with identical inputs and method, changing only histogram binning.
\end{itemize}
\centering
\begin{tabular}{lccc}
\toprule
Setup & Recommended $N_{\mathrm{iter}}$ & Toy closure mean & Refold $\chi^2/\mathrm{ndf}$ \\
\midrule
50 bins (current nominal) & 7 & 0.983 & 2.278 \\
25 bins (cross-check) & 4 & 0.978 & 2.421 \\
\bottomrule
\end{tabular}
\vspace{0.6em}
\begin{itemize}
  \item Conclusion: the previous 4--5 result is consistent with a coarser binning setup.
  \item Current 50-bin setup needs more iterations to reach comparable stability.
\end{itemize}
\end{frame}

\begin{frame}{Input and Response Distributions}
\begin{columns}[T]
\column{0.5\textwidth}
\centering
\includegraphics[width=\linewidth,height=0.55\textheight,keepaspectratio]{../results/dist_inputs_data_mcreco_mcgen.pdf}
\column{0.5\textwidth}
\centering
\includegraphics[width=\linewidth,height=0.55\textheight,keepaspectratio]{../results/dist_response_counts_matrix.pdf}
\end{columns}
\vspace{0.1em}
\begin{itemize}
  \item Left: data reco, MC reco, MC gen truth input spectra.
  \item Right: reco-vs-gen response matrix used by unfolding.
  \item All 1D thrust spectra are shown as differential distributions: $1/N\,dN/dT$ (log y-axis).
\end{itemize}
\end{frame}

\begin{frame}{Unfolded and Refolded Distributions}
\begin{columns}[T]
\column{0.5\textwidth}
\centering
\includegraphics[width=\linewidth,height=0.55\textheight,keepaspectratio]{../results/dist_unfolded_selected_iterations.pdf}
\column{0.5\textwidth}
\centering
\includegraphics[width=\linewidth,height=0.55\textheight,keepaspectratio]{../results/dist_refolded_vs_data_selected_iterations.pdf}
\end{columns}
\vspace{0.1em}
\begin{itemize}
  \item Left: unfolded distributions for $N=4,7,10$ compared to MC truth.
  \item Right: corresponding refolded spectra compared to data reco.
\end{itemize}
\end{frame}

\begin{frame}{Shape Comparison of Unfolded Distributions}
\centering
\includegraphics[width=0.97\linewidth,height=0.72\textheight,keepaspectratio]{../results/dist_unfolded_shape_selected_iterations.pdf}
\vspace{0.1em}
\begin{itemize}
  \item Shape-normalized view highlights where higher iterations start to shift spectral shape.
  \item $N=7$ remains close to truth shape while avoiding larger fluctuations seen at higher iteration.
\end{itemize}
\end{frame}

\begin{frame}{Main Trends}
\begin{columns}[T]
\column{0.5\textwidth}
\centering
\includegraphics[width=\linewidth,height=0.53\textheight,keepaspectratio]{../results/chi2ndf_toy_closure_counts_vs_iterations.pdf}
\column{0.5\textwidth}
\centering
\includegraphics[width=\linewidth,height=0.53\textheight,keepaspectratio]{../results/chi2ndf_refold_shape_vs_iterations.pdf}
\end{columns}
\vspace{0.1em}
\begin{itemize}
  \item Toy closure rises with iteration (noise amplification at high iterations).
  \item Refold agreement with data improves with iteration.
  \item A compromise point is required rather than pure minimization of one metric.
\end{itemize}
\end{frame}

\begin{frame}{Convergence Check: Consecutive Iterations}
\begin{columns}[T]
\column{0.5\textwidth}
\centering
\includegraphics[width=\linewidth,height=0.53\textheight,keepaspectratio]{../results/chi2ndf_iter_delta_counts_vs_iterations.pdf}
\column{0.5\textwidth}
\centering
\includegraphics[width=\linewidth,height=0.53\textheight,keepaspectratio]{../results/chi2ndf_iter_delta_diffxsec_vs_iterations.pdf}
\end{columns}
\vspace{0.1em}
\begin{itemize}
  \item Left: $\chi^2/\mathrm{ndf}(\mathrm{unfolded}_{N},\mathrm{unfolded}_{N-1})$ in counts.
  \item Right: $\chi^2/\mathrm{ndf}$ between consecutive differential spectra $(1/N\,dN/dT)_N$ and $(1/N\,dN/dT)_{N-1}$.
  \item It decreases monotonically: each extra iteration changes the result less than the previous step.
  \item Differential metric drops from 0.231 ($N=4$) to 0.121 ($N=5$), 0.071 ($N=6$), 0.044 ($N=7$).
  \item This places the $\chi^2/\mathrm{ndf}\sim 0.05$ convergence scale around $N\approx 7$ for the 50-bin setup.
\end{itemize}
\end{frame}

\begin{frame}{Overfitting Study: Statistical Growth}
\begin{columns}[T]
\column{0.5\textwidth}
\centering
\includegraphics[width=\linewidth,height=0.53\textheight,keepaspectratio]{../results/overfit_avg_frac_unc_vs_iterations.pdf}
\column{0.5\textwidth}
\centering
\includegraphics[width=\linewidth,height=0.53\textheight,keepaspectratio]{../results/overfit_avg_pull_width_vs_iterations.pdf}
\end{columns}
\vspace{0.1em}
\begin{itemize}
  \item Left: average toy fractional uncertainty $\langle \sigma_{\mathrm{toy}}/\mathrm{truth} \rangle$ grows with iteration.
  \item Right: fluctuation width $\langle \sigma_{\mathrm{toy}}/\sqrt{\mathrm{truth}} \rangle$ crosses $\sim 1$ around $N\approx 8$ and keeps increasing.
  \item Interpretation: higher iterations amplify statistical noise, consistent with IBU overfitting behavior.
\end{itemize}
\end{frame}

\begin{frame}{Overfitting Study: Correlations and MSE}
\begin{columns}[T]
\column{0.5\textwidth}
\centering
\includegraphics[width=\linewidth,height=0.53\textheight,keepaspectratio]{../results/overfit_bias_variance_mse_vs_iterations.pdf}
\column{0.5\textwidth}
\centering
\includegraphics[width=\linewidth,height=0.53\textheight,keepaspectratio]{../results/overfit_mean_abs_offdiag_corr_vs_iterations.pdf}
\end{columns}
\vspace{0.1em}
\begin{itemize}
  \item Variance/MSE increase strongly with iteration; this is the primary overfitting indicator here.
  \item Mean absolute off-diagonal correlation stays at a similar level and is less discriminating than variance growth.
  \item Practical takeaway: avoid very high iterations ($\gtrsim 10$) where statistical inflation is clear.
\end{itemize}
\end{frame}

\begin{frame}{Representative Values}
\centering
\begin{tabular}{rcccc}
\toprule
$N_{\mathrm{iter}}$ & Refold $\chi^2/\mathrm{ndf}$ & Toy closure mean & Toy RMS & $\Delta_{N,N-1}$ \\
\midrule
1  & 6.948 & 0.173 & 0.074 & -- \\
4  & 3.119 & 0.585 & 0.197 & 0.320 \\
5  & 2.731 & 0.721 & 0.260 & 0.121 \\
6  & 2.466 & 0.828 & 0.236 & 0.137 \\
7  & 2.278 & 0.936 & 0.263 & 0.097 \\
8  & 2.141 & 1.113 & 0.344 & 0.070 \\
10 & 1.958 & 1.156 & 0.377 & 0.040 \\
15 & 1.739 & 1.871 & 0.706 & 0.014 \\
20 & 1.641 & 2.401 & 0.838 & 0.007 \\
\bottomrule
\end{tabular}

\vspace{0.8em}
Selection rule used: require refold $\chi^2/\mathrm{ndf} \lesssim 2.5$, toy closure near 1, and check $\Delta_{N,N-1}$ for convergence.
\end{frame}

\begin{frame}{Recommendation and Takeaways}
\begin{block}{Recommended iteration}
\centering
\Large \textbf{$N_{\mathrm{iter}}=5$}
\end{block}

\begin{itemize}
  \item At $N_{\mathrm{iter}}=5$: refold $\chi^2/\mathrm{ndf}=2.731$, toy closure mean $=0.721$.
  \item Differential consecutive-change metric at $N=5$ is already in the low regime ($\Delta_{N,N-1}=0.121$), with a further drop at $N=6$.
  \item Overfit study: statistical fluctuations increase with iteration; $N=5$ keeps uncertainties lower than higher-iteration choices.
  \item Lower iterations underfit data (high refold $\chi^2/\mathrm{ndf}$).
  \item Higher iterations ($\gtrsim 10$) increase toy closure $\chi^2/\mathrm{ndf}$, consistent with over-unfolding.
  \item Practical working point: use 5 iterations nominal and evaluate $\pm1$ iterations for regularization uncertainty.
\end{itemize}
\end{frame}

\begin{frame}{Final Analysis Choice}
\begin{block}{Adopted Iteration Prescription}
\centering
\Large Nominal: \textbf{$N_{\mathrm{iter}}=5$} \\
\Large Regularization systematics: \textbf{$N_{\mathrm{iter}}=4,6$}
\end{block}
\begin{itemize}
  \item Motivation: convergence behavior is already flattening by $N=5$ while statistical growth remains controlled.
  \item Uncertainty treatment: take the up/down envelope from 4 and 6 iterations.
  \item This will be propagated as unfolding regularization systematic in final results.
\end{itemize}
\end{frame}

\begin{frame}{Proposed Next Steps}
\begin{enumerate}
  \item Repeat the same scan for key detector/systematic variations (NES up/down, NER, NTPC7, PT04).
  \item Run the nominal regularization envelope with 4, 5, and 6 iterations for final uncertainty extraction.
  \item Add bin-by-bin pull distributions from toys to identify bins where over-regularization/overfitting appears first.
  \item Propagate iteration uncertainty as a regularization systematic using the 4/6 variation around nominal 5.
  \item Cross-check with an alternative regularization strategy (e.g. SVD unfolding) as a robustness test.
\end{enumerate}
\end{frame}

\end{document}
